\documentclass{ximera}

\input{../../preamble.tex}

\title{Exponential Analysis}

\begin{document}

\begin{abstract}
%Stuff can go here later if we want!
\end{abstract}
\maketitle




Precalculus is largely about the structure of the Elementary Functions and how we reason about them.  This includes power functions; radical and root functions; linear, quadratic, and other polynomial functions; rational functions; exponential functions; logarithmic functions; and a library of trigonometric functions (which are introduced in this course).


Our focus is analyzing these functions.





\begin{itemize}
     \item \textbf{\textcolor{red!70!black}{Domain}} 
     \item \textbf{\textcolor{red!70!black}{Zeros}} 
     \item \textbf{\textcolor{red!70!black}{Continuity}} \\
     \text{    }$\circ$ \textbf{\textcolor{purple!85!blue}{discontinuities}}\\
     \text{    }$\circ$ \textbf{\textcolor{purple!85!blue}{singularities}}\\
     \item \textbf{\textcolor{red!70!black}{End-Behavior}} 
     \item \textbf{\textcolor{red!70!black}{Behavior}} \\
     \text{    }$\circ$ \textbf{\textcolor{purple!85!blue}{intervals where increasing}}\\
     \text{    }$\circ$ \textbf{\textcolor{purple!85!blue}{intervals where decreasing}}\\
     \item \textbf{\textcolor{red!70!black}{Global Maximum and Minimum}} 
     \item \textbf{\textcolor{red!70!black}{Local Maximums and Minimums}} 
     \item \textbf{\textcolor{red!70!black}{Range}} 
     \item \textbf{\textcolor{blue!55!black}{...and we would like a nice graph}} 
\end{itemize}






As a segue from the first half of Precalculus, let's quickly review our growing library of Elementary Functions, which are linear combinations and compositions of the Core Functions with linear functions.


\subsection*{Core Functions} 

The Core Functions are our basic building blocks. The Elementary Functions will grow from of the Core Functions.

\begin{itemize}
\item{Constant Functions: $C(x) = c$}
\item{The Identity Function: $Id(x) = x$}
\item{Power Functions: $P(x) = x^n$}
\item{Root Functions: $R(x) = \sqrt[n]{x}$}
\item{Exponential Functions: $E(x) = r^x$}
\item{Logarithmic Functions: $L(x) = \log_r(x)$}
\item{Sine and Cosine: $S(x) = \sin(x)$ and $C(x) = \cos(x)$}
\item{The Absolute Value Function: $AV(x) = |x|$}
\end{itemize}


\subsection*{Elementary Functions} 

By forming sums, differences, products, quotients, compositions, and inverses of these Core Functions, we obtain the \textbf{Elementary Functions}.


\begin{itemize}
\item{Constant}
\item{Linear}
\item{Quadratic}
\item{Polynomial}
\item{Rational}
\item{Radical/Root}
\item{Exponential}
\item{Shifted Exponential}
\item{Logarithmic}
\item{Absolute Value}
\item{Trigonometric}
\item{Inverse Trigonometric}
\end{itemize}




After a quick reminder of the Core Functions characteristics and features, we'll begin analyzing exponential functions in this section.   





$\blacktriangleright$ \textbf{\textcolor{blue!55!black}{Basic Exponential Functions}} 


The template for the basic exponential function looks like \textbf{\textcolor{blue!55!black}{$exp(x) = A \cdot r^x$}} or  (\textbf{\textcolor{blue!55!black}{$A \cdot r^{B \, x + C}$}}). 

\begin{itemize}
\item $A$ is the \textbf{\textcolor{purple!85!blue}{leading coefficient}}.  It primarily determines the sign of the function values.
\item $r$ is the \textbf{\textcolor{purple!85!blue}{base}}. It determines how the function values grow, depending on whether $0 < r < 1$ or $1 < r$.
\end{itemize}





$\blacktriangleright$ \textbf{\textcolor{blue!55!black}{Shifted Exponential Functions}}   

The template for a shifted exponential function looks like \textbf{\textcolor{blue!55!black}{$shexp(x) = A \cdot r^x + D$}}  or  (\textbf{\textcolor{blue!55!black}{$A \cdot r^{B \, x + C} + D$}}). 

\begin{itemize}
\item $C$ describes a \textbf{\textcolor{purple!85!blue}{shift}} in the domain from the core exponential function.  
\item $D$ describes a \textbf{\textcolor{purple!85!blue}{shift}} in the function value from the core exponential function. 
\end{itemize}

The first term in the formula for a shifted exponential function is, by itself, an exponential function:

\[
a \cdot r^{b \, x + c} = a \cdot r^{b \, x} \cdot  r^c = (a \cdot r^c) \cdot r^{b \, x} = (a \cdot r^c) \cdot (r^b)^x 
\]

It \textbf{\textcolor{purple!85!blue}{CAN}} be written in the form $A \cdot R^x$, where $A = a \cdot r^c$ and $R = r^b$.


Algebraically, the constant term, ``+ D'', cannot be absorbed into the exponential form.  However, shifting functions is a very common operation and so, we'll include $A \cdot r^{B \, x + C} + D$ in our general ideas for exponential functions.





\subsection*{Learning Outcomes}


\begin{sectionOutcomes}
In this section, students will 

\begin{itemize}
\item describe the full picture of all exponential and shifted exponential functions and their graphs.
\end{itemize}
\end{sectionOutcomes}












\begin{center}
\textbf{\textcolor{green!50!black}{ooooo-=-=-=-ooOoo-=-=-=-ooooo}} \\

more examples can be found by following this link\\ \link[More Examples of Exponential Functions]{https://ximera.osu.edu/csccmathematics/precalculus2/precalculus2/expFunctions/examples/exampleList}

\end{center}









\end{document}
